\documentclass[11pt,a4paper]{article}

% =============================================================================
% Provisional regression-output companion -- reform-window DDD
% =============================================================================
% Standalone document holding the current regression outputs from the reform-
% window DDD design (DA15 + IDA15). The design is NOT pinned down: the
% identification choices (heterogeneity vs treatment/control, covariate set,
% session-choice for IDA, etc.) are still being iterated, and the user has
% flagged that the descriptive foundations must come first. The numbers here
% are the current state of play; they may be revised or replaced once the
% descriptive companion (descriptive_facts.tex) settles the design language.
%
% Content was migrated 2026-05-17 from thesis/paper/provisional.tex section 1.
% The descriptive bidding-pattern content from that file lives in the sibling
% document thesis/provisional/descriptive_facts.tex.

\PassOptionsToPackage{table}{xcolor}
\usepackage{amsmath, amssymb, amsthm, mathtools}
\usepackage{graphicx, booktabs, subcaption, tikz, float}
\graphicspath{{../../figures/thesis/}}
\usepackage{array, makecell, multirow}
\usepackage{xcolor}
\usepackage{longtable}
\usepackage[hidelinks]{hyperref}
\usepackage{geometry}
\geometry{a4paper, margin=2.5cm}
\usepackage{caption}
\captionsetup{font=small, labelfont=bf}

\usepackage{url}
\urlstyle{tt}
\Urlmuskip=0mu plus 1mu\relax
\def\UrlBreaks{\do\/\do\_\do\-\do\.\do\:\do\,}
\newcommand{\code}[1]{\nolinkurl{#1}}

\title{Provisional: regression outputs --- reform-window DDD (DA15 + IDA15)\\[3pt]\large MTU15 thesis companion (design NOT pinned down)}
\author{Pablo Paramio}
\date{\today}

\begin{document}
\maketitle

\noindent\emph{\textbf{Status.} This file collects the current reform-window DDD outputs. The identification design is being iterated --- the descriptive foundations live in the sibling file \code{thesis/provisional/descriptive_facts.tex} and the framing of "heterogeneity vs treatment/control", covariate choice, and session-choice for the IDA outcomes is still in flux. Read the numbers here as the current state of play, not as final results.}

\tableofcontents

\bigskip

% =====================================================================
\section{Reform-window identification (DA15 + IDA15)}\label{sec:prov:reform_window_ddd}
% =====================================================================

\noindent The headline thesis question is the effect of \textbf{MTU15 granularity} on strategic bidding. Reforzada (the post-blackout RT2 pay-as-bid + zonal-monopoly windfall channel documented in the descriptive companion, Part B) is a parallel \emph{confounding} mechanism --- analogous in spirit to the operational/solar midday channel in the main paper (\code{paper.tex}~\S2.4): a real economic phenomenon that we want to characterise but \emph{exclude from the identifying contrast} so as not to mix channels. This section's design picks reform windows in which reforzada is held constant, then differences out residual seasonality via a 2024 placebo.

\paragraph{Design.} Two reform-windows, each chosen so that reforzada is constant within both the pre- and post-treatment sub-windows:

\begin{table}[H]
\centering
\caption{Reform-window design. Reforzada is the post-blackout RT2 pay-as-bid regime (active 2025-04-28 onward). The 2024 placebo uses calendar-matched dates. PT assumption: YoY seasonal change in (crit $-$ flat) was the same in 2024 as in 2025.}
\label{tab:reform_windows}
\footnotesize
\setlength{\tabcolsep}{4pt}
\begin{tabular}{@{}l l l c c@{}}
\toprule
Reform / sample & Pre-window & Post-window & Reforzada & Length pre/post \\
\midrule
\textbf{IDA15 (2025)} & 2024-12-09 -- 2025-03-18 & 2025-03-19 -- 2025-04-27 & off & 100 / 40 d \\
IDA15 placebo (2024)  & 2023-12-09 -- 2024-03-18 & 2024-03-19 -- 2024-04-27 & off & matching \\
\addlinespace
\textbf{DA15 (2025)} & 2025-04-28 -- 2025-09-30 & 2025-10-01 -- 2026-02-13 & on  & 155 / 135 d \\
DA15 placebo (2024)  & 2024-04-28 -- 2024-09-30 & 2024-10-01 -- 2025-02-13 & off & matching \\
\bottomrule
\end{tabular}
\end{table}

\noindent The temporal design holds reforzada constant within each window (on for DA15, off for IDA15), and the 2024 placebo absorbs YoY seasonality that the within-day DiD alone could not. Each pre-window is $\ge 3$ months; DA15 has 5 months post-reform.

\begin{figure}[H]
\centering
\includegraphics[width=0.95\textwidth]{../../figures/working/fig_reform_windows_timeline.pdf}
\caption{\textbf{Timeline of reforms and DDD windows.} Vertical guides: ISP15 (2024-12-09), MTU15-IDA (2025-03-19), Iberian blackout (2025-04-28), MTU15-DA (2025-10-01). Shaded grey: reforzada zone. Blue bars: pre-reform sub-windows; red bars: post-reform. 2024 placebo rows shown with lighter shading.}
\end{figure}

\paragraph{Triple-difference specification.}
Let $Y_{f,u,d,h}$ be the outcome at firm $f$, unit $u$, date $d$, clock-hour $h$. Define $\text{crit}_h = 1$ if $h \in \{5\!-\!9,\,16\!-\!22\}$ (treated hour-class), restrict sample to $\{\text{crit}, \text{flat}\}$. Let $\text{post}_d = 1$ if $d$ is after the reform date, $\text{y25}_d = 1$ if $d$ falls in the 2025 (treatment) window. Then:
\begin{multline}
  Y_{f,u,d,h} = \alpha + \beta_1 \text{crit}_h + \beta_2 \text{post}_d + \beta_3 \text{y25}_d
    + \beta_4 (\text{crit} \times \text{post}) + \beta_5 (\text{crit} \times \text{y25}) \\
    + \beta_6 (\text{post} \times \text{y25})
    + \boxed{\beta_7 (\text{crit} \times \text{post} \times \text{y25})}
    + \gamma_f + \mu_u + \delta_{\text{DOW}(d)} + \varepsilon_{f,u,d,h}.
\end{multline}
$\beta_7$ is the headline coefficient: the change in the (critical $-$ flat) differential at the reform date, \emph{net of} the corresponding 2024 seasonal change. FE: firm $\gamma_f$ + unit $\mu_u$ + day-of-week $\delta_{\text{DOW}}$. SE clustered at the date level. Equivalently $\beta_7 = \Delta\Delta Y^{2025} - \Delta\Delta Y^{2024}$ where $\Delta\Delta Y = (\bar Y_{\text{post,crit}} - \bar Y_{\text{post,flat}}) - (\bar Y_{\text{pre,crit}} - \bar Y_{\text{pre,flat}})$.

\paragraph{Event-study spec (visual pre-trends).} For each outcome, plot the monthly differential separately for the 2025 treatment year and the 2024 placebo year, anchored to $\tau =$ months from the reform date:
\[
  \Delta Y_{\tau,y} \;=\; \bar Y_{\text{crit}}(\tau, y) - \bar Y_{\text{flat}}(\tau, y), \qquad y \in \{2024, 2025\}.
\]
Identifying assumption test: for $\tau < 0$ the two series should overlap (parallel pre-trends); at $\tau \ge 0$ the 2025 series diverges from the 2024 series by exactly the treatment effect.

\paragraph{Identifying assumptions.}
\textit{Formal PT (triple-difference).}
\[
  \mathbb{E}\big[\Delta\Delta Y(0) \mid \text{y25}=1\big] \;=\; \mathbb{E}\big[\Delta\Delta Y(0) \mid \text{y25}=0\big],
\]
where $\Delta\Delta Y(0)$ is the untreated within-day pre/post change in the critical-flat differential.

\textit{Threats killed by the design.}
\begin{itemize}
   \item \emph{Seasonality of the critical-flat spread} (summer-to-winter heating/AC patterns): differenced out by the 2024 placebo on the same calendar months.
   \item \emph{Day-level cost shocks} (TTF gas, CO$_2$, demand spikes): within-day differencing on each date absorbs anything that hits both hour-classes proportionally.
   \item \emph{Asymmetric gas-cost confound}: ruled out empirically in \code{paper.tex}~\S2.4 --- CCGT price-setting share 9.1\% crit vs 9.0\% flat, so TTF shocks cannot propagate asymmetrically.
   \item \emph{Operational midday channel} (solar drives prices down at midday): midday hours $\in \{11\!-\!14\}$ excluded from the control set; only flat hours $\in \{1,2,3\}$ enter as control.
   \item \emph{Reforzada / RT2 pay-as-bid windfall}: held constant within each window by design (the central insight). Characterised as a confounder in the descriptive companion, Part B (REE post-clearing CCGT inclusion, geographic CCGT concentration, Cournot reading $+$ CNMC sanctions).
\end{itemize}

\textit{Threats remaining.}
\begin{itemize}
   \item \emph{Within-firm portfolio composition shifts}: addressed by unit FE; sensitivity check restricts to units active in both years.
   \item \emph{Reform-anticipation effects}: tested by the event-study leads $\beta_\tau$ for $\tau \in \{-1, -2\}$.
\end{itemize}

\paragraph{What this DDD identifies (and what it does not).}
The within-day differencing on (crit $-$ flat) absorbs day-level common shocks; the 2024 placebo absorbs YoY seasonal variation in the spread itself. Whether the identified parameter is the \emph{differential} reform effect or the \emph{absolute} effect on critical hours depends on whether flat hours are exposed to MTU15. This is a market-specific empirical question; the existing project evidence (descriptive companion, $D_w$ decomposition section, DA Oct--Dec 2025) shows:
\begin{itemize}
    \item \textbf{DA market}: flat hours show essentially zero within-hour granularity activity (GN, GE, HC: 0\%--0.05\% on both $\Delta p$ and $\Delta q$ channels; IB: 0.89\%/1.03\%, an order of magnitude below its critical-hour shaping). Critical hours, by contrast, show GN at 18.4\% $\Delta q$-active. Flat hours are approximately MTU15-untreated for CCGT bidding. \emph{The DA-cell DDDs on prices and on $q_1$ should therefore be read as close to the absolute critical-hour effect, not a within-day differential between two treated arms.}
    \item \textbf{IDA market} (pre-blackout, see descriptive companion): flat-hour cell activity is 7--21\% across pivotal firms, against critical-hour 9--30\%. Both hour-classes show within-hour activity in the IDA. \emph{The IDA-cell DDD identifies a true differential, not an absolute effect on critical hours.}
\end{itemize}
Same logic for pivotal-vs-non-pivotal in the $D_w$ decomposition: the table only breaks out the four pivotal firms because non-pivotal CCGT participation in the relevant zones is negligible. For granularity-driven outcomes the non-pivotal-as-control reading is more defensible than for aggregate outcomes like $q_2$, where the older 2-way DiD in \code{paper.tex} appendix shows non-pivotal firms do respond (with sign reversal). The framing is outcome- and market-specific: closer to treatment/control for DA-granularity-shaping outcomes, closer to heterogeneity for IDA and for aggregate $q_2$.

\paragraph{Heterogeneity decomposition by pivotal firm status.}
For firm-level outcomes ($q_1$, $q_2$) we report the DDD twice: pooled over the CCGT sample, and decomposed by pivotal status (Big-4: IB, GE, GN, HC, EDP-PT; non-pivotal: Repsol, Engie, TotalEnergies, Moeve). The decomposition is a \emph{heterogeneity analysis}, not a fourth difference: non-pivotal CCGT firms are exposed to the same MTU15 reform as pivotal firms and the older 2-way DiD in \code{paper.tex} appendix in fact shows them responding (the $q_2$ sign on (crit $\times$ post) reverses between pivotal and non-pivotal subsamples). We do not assume they are an "untreated control"; we read the contrast as how the reform's effect varies with market power. Operationally this is implemented as a saturated OLS-FE regression with all crit/post/y25/pivotal interactions (the pivotal main effect is firm-invariant and absorbed by $\gamma_f$), yielding a coefficient on (crit $\times$ post $\times$ y25 $\times$ piv) that is mathematically the difference of the two sub-sample DDDs. We label it $\delta_{\text{piv}}$ rather than $\beta_{1234}$ to make the heterogeneity-not-fourth-difference reading explicit:
\[
  \delta_{\text{piv}} \;\equiv\; \widehat{\text{DDD}}_{\text{pivotal}} \;-\; \widehat{\text{DDD}}_{\text{non-pivotal}}.
\]
SE clustered at the firm level (pivotal does not vary within firm; $G \approx 9$ --- wild-cluster bootstrap as a robustness step is flagged below). The DDDD-style name and the "quadruple difference" framing previously used in this section have been retired: in light of the credibility-quadratic-in-differences argument (Roth, paraphrased in Borusyak's lecture notes), and given the empirical evidence that non-pivotal firms are not a clean control, the heterogeneity-decomposition interpretation is the only one we sustain.

\paragraph{Outcomes and joint $q_1 / q_2$ sign-quadrants.}
Two firm-level quantity outcomes and two market-level price outcomes. The firm-level pair is informative \emph{only when read jointly}: $q_1$ (DA-cleared MWh) and $q_2$ (signed intraday MWh) jointly characterise the firm's strategic posture, with each $(\text{sign}(\beta_{q_1}), \text{sign}(\beta_{q_2}))$ quadrant mapping to a distinct economic narrative. Table~\ref{tab:reform_outcomes} lays this out. $D_w$, $\Delta p$, $\Delta q$ are \emph{not} DDD outcomes because they require within-hour quarters that do not exist pre-MTU15; they appear as post-only cross-section in the descriptive companion ($D_w$ decomposition and $D_w$ distribution sections).

\begin{table}[H]
\centering\footnotesize
\setlength{\tabcolsep}{6pt}
\caption{Joint $(q_1, q_2)$ sign quadrants and the firm-behaviour narrative each implies. The DDD coefficient on critical-flat-spread for each quantity outcome gives the row/column sign. Each cell is a coherent story about how dispatchable firms responded to the reform; only one is realised in the data.}
\label{tab:reform_outcomes}
\begin{tabular}{@{}p{2.6cm} p{5.5cm} p{5.5cm}@{}}
\toprule
 & $\beta(q_1) > 0$ (more DA-cleared) & $\beta(q_1) < 0$ (less DA-cleared) \\
\midrule
$\beta(q_2) > 0$ (more IDA selling) & \emph{Production expansion}: firms scale up across both stages (unlikely under unchanged demand). & \emph{Ito--Reguant withholding}: classic strategic relabelling --- withhold in DA to drive up clearing, sell volume back in IDA. \\
\addlinespace
$\beta(q_2) < 0$ (less IDA selling) & \emph{Operational refinement}: granular DA absorbs within-hour variation that previously needed IDA repositioning; auctions clear better, less intraday correction needed. & \emph{Broader disengagement}: firms reduce participation in \emph{both} auction-cleared stages, with the outside option being REE post-clearing redispatch (RT2 pay-as-bid). Not Ito--Reguant. \\
\bottomrule
\end{tabular}
\end{table}

\noindent For prices, the within-day spread response separately tests the operational-refinement vs strategic-exploitation directions: negative ($\beta_7 < 0$, spread compresses) if granularity lets non-pivotal competitors fill critical-hour 15-min slots; positive ($\beta_7 > 0$, spread widens) if pivotal firms exploit the new pricing dimension.

\paragraph{How to read Figure~\ref{fig:event_study}.}
Each panel is a single (outcome $\times$ reform) cell. The horizontal axis is $\tau =$ months relative to the reform date (e.g.\ $\tau = 0$ is March 2025 for IDA15 and October 2025 for DA15). The vertical axis is the monthly mean of the within-day differential $\bar Y_{\text{crit}}(\tau) - \bar Y_{\text{flat}}(\tau)$ in that month --- a scalar that tells us how much higher the critical-hour outcome is relative to the same day's flat-hour outcome.

Each panel has two coloured lines:
\begin{itemize}
\item \textcolor{red}{Red (2025)}: the actual treatment year. At $\tau = 0$, MTU15 is implemented.
\item \textcolor{blue}{Blue (2024)}: the calendar-matched placebo. At $\tau = 0$, the analogous calendar month one year earlier --- no reform, just seasonality.
\end{itemize}

The DDD coefficient $\beta_7$ is a difference of \emph{changes} (not levels). In raw form the two lines can sit at different levels (e.g., 2025 prices systematically higher than 2024 because of inflation) and the design still works --- the level shift is absorbed by the $\text{y25}$ main effect $\beta_3$ and what matters is that the (red $-$ blue) gap stays constant pre-reform and shifts post-reform.

The figure normalises each line by subtracting its value at $\tau = -1$, so both start at zero at the omitted reference month. After normalisation the visual test is the standard event-study form:
\begin{itemize}
\item For $\tau < 0$ (pre-reform): red and blue should \textbf{overlap} (parallel-trends test).
\item At $\tau \geq 0$ (post-reform): a persistent vertical separation between red and blue is the treatment effect.
\item $\beta_7$ is the average post-period vertical separation between the two normalised lines.
\end{itemize}

\begin{figure}[H]
\centering
\includegraphics[width=\textwidth]{../../figures/working/fig_event_study_ddd.pdf}
\caption{Event-study (reduced-form, normalised). Rows = outcomes (DA price, IDA price, $q_1$, $q_2$); columns = reforms (IDA15, DA15). Each panel plots the monthly (critical $-$ flat) within-day differential, anchored to $\tau =$ months from the reform date and \emph{normalised by subtracting the $\tau = -1$ value}, so both lines start at 0 at $\tau = -1$. Red = 2025 (treatment); blue = 2024 (placebo). Pre/post regions lightly shaded. Reduced-form monthly means; full regression-based event-study with CI follows the headline DDD.}
\label{fig:event_study}
\end{figure}

\textit{Pre-trend reading.} DA15: pre-period overlap is reasonable for prices; $q_1$ shows the 2025 series sitting below the 2024 series \emph{already before $\tau = 0$}, signalling that the reforzada channel pre-bends the differential before MTU15-DA fires --- a partial PT violation for $q_1$ on this window. IDA15: pre-period overlap looks reasonable for prices; post-period is too short to read visually. The headline DDD with date-clustered SE follows.

\paragraph{Estimators in practice: OLS-FE and doubly-robust DiD on prices in levels.}
We report two estimators for the same DDD parameter, each with three to four specifications. Both use clearing prices in EUR/MWh as the outcome (in levels), and both target the same coefficient: the (critical $-$ flat) differential response to the reform, net of the 2024 calendar placebo.

\textit{Panel construction.} For each (reform window, market) combination $w \in \{\text{DA15}, \text{IDA15}\} \times \{\text{DA}, \text{IDA-S1}\}$ we build two analytic panels:
\begin{itemize}
\item \emph{Period panel}: one row per (date, ISP-period) tuple. Pre-MTU15 days contribute 1 observation per clock-hour (MTU60); post-MTU15 days contribute 4 (MTU15). Sample restricted to clock-hours in \{critical $\cup$ flat\}. Used for the OLS-FE estimator. Sizes: 13\,888 obs (DA15-DA), 19\,138 obs (DA15-IDA), 3\,934 obs (IDA15-DA), 5\,530 obs (IDA15-IDA).
\item \emph{Spread panel}: one row per date with the within-day means $\bar Y^{\text{crit}}_d = \frac{1}{|H_{\text{crit}}|}\sum_{p \in H_{\text{crit}}} Y_{d,p}$ and $\bar Y^{\text{flat}}_d$ as two columns. Used for the DR-DiD estimator (one drdid call per hour-class). Sizes: 584 dates (DA15-DA), 554 (DA15-IDA), 281 (IDA15-DA), 278 (IDA15-IDA).
\end{itemize}

\textit{OLS-FE DDD.} On the period panel, three specifications, each with the same DDD coefficient $\beta_7$ but different fixed-effect structures:
\begin{align*}
\text{(M1) sparse:}\quad & Y_{d,p} = \alpha + \sum_{k=1}^{7} \beta_k Z_k + \gamma_h + \delta_{\text{DOW}(d)} + \varepsilon_{d,p}, \\
\text{(M2) month FE:}\quad & Y_{d,p} = \alpha + \sum_{k=1}^{7} \beta_k Z_k + \gamma_h + \delta_{\text{DOW}(d)} + \eta_{\text{month}(d)} + \varepsilon_{d,p}, \\
\text{(M3) date FE:}\quad & Y_{d,p} = \alpha + \beta_1 \text{crit} + \beta_4 (\text{crit}\!\times\!\text{post}) + \beta_5 (\text{crit}\!\times\!\text{y25}) \\
& \qquad\qquad\quad + \beta_7 (\text{crit}\!\times\!\text{post}\!\times\!\text{y25}) + \gamma_h + \delta_{\text{DOW}(d)} + \nu_d + \varepsilon_{d,p}.
\end{align*}
Here $(Z_1,\ldots,Z_7) = (\text{crit},\,\text{post},\,\text{y25},\,\text{crit}\!\times\!\text{post},\,\text{crit}\!\times\!\text{y25},\,\text{post}\!\times\!\text{y25},\,\text{crit}\!\times\!\text{post}\!\times\!\text{y25})$, $\gamma_h$ is a clock-hour FE, $\delta_{\text{DOW}(d)}$ is a day-of-week FE, $\eta_{\text{month}(d)}$ is a calendar-month FE, and $\nu_d$ is a date FE that absorbs all date-level terms in M3 (so the date-level $\beta_2,\beta_3,\beta_6$ drop out and $\beta_7$ is identified purely from within-date variation). SE clustered at the date level via \texttt{reghdfe ..., vce(cluster date)}.

\textit{Doubly-robust DiD (Sant'Anna--Zhao 2020).} On the spread panel, two 2$\times$2 doubly-robust DiDs per (window, market), one on each hour-class outcome:
\[
\widehat{\text{ATT}}^{\text{DR}}_{c} \;=\; \mathbb{E}_n\!\left[\Big(\,w_{1}(D)\,-\,w_{0}(D,X;\hat p)\Big)\,\Big(\,(Y^{c}_{\text{post}}-Y^{c}_{\text{pre}})\,-\,\big(\hat m^{0}_{\text{post}}(X)-\hat m^{0}_{\text{pre}}(X)\big)\,\Big)\right],
\]
with $c \in \{\text{crit}, \text{flat}\}$, treatment indicator $D = \text{y25}_d$, time indicator $t = \text{post}_d$, covariates $X = (\text{i.dow}, \text{i.month})$, and
\[
w_{1}(D) = \frac{D}{\mathbb{E}_n[D]}, \qquad
w_{0}(D, X; \hat p) = \frac{\hat p(X)(1-D)/(1-\hat p(X))}{\mathbb{E}_n[\hat p(X)(1-D)/(1-\hat p(X))]}.
\]
The default \code{drdid} estimator (\code{drimp}) uses inverse-probability-tilting for $\hat p(X)$ and weighted least squares for the outcome model $\hat m^{0}(X)$. Three nuisance-model specifications are reported:
\begin{itemize}
\item \emph{DR1}: no covariates ($X = \emptyset$). Reduces to the unconditional 2$\times$2 DiD on the spread $Y^{c}$.
\item \emph{DR2}: $X = (\text{i.dow}, \text{i.month})$, asymptotic SE via the influence function.
\item \emph{DR3}: $X = (\text{i.dow}, \text{i.month})$, wild-bootstrap inference (999 Mammen draws, \texttt{drdid ..., wboot rseed(4242)}).
\end{itemize}

\textit{DDD via the difference of two doubly-robust DiDs (the headline).} The triple difference is the difference of the two hour-class ATTs:
\[
\widehat{\text{DDD}} \;=\; \widehat{\text{ATT}}^{\text{DR}}_{\text{crit}} \,-\, \widehat{\text{ATT}}^{\text{DR}}_{\text{flat}}.
\]
Joint inference uses the recentered influence functions $\hat\psi^{c}_i$ saved by drdid's \texttt{stub(}$c$\texttt{)} option, which satisfy $\sqrt N(\widehat{\text{ATT}}^{c} - \text{ATT}^{c}) = N^{-1/2}\sum_i \hat\psi^{c}_i + o_p(1)$ (Sant'Anna--Zhao 2020, Theorem 1). Because the crit and flat means are computed on the same dates (within-date dependence), the variance of the DDD is taken from the variance of the IF \emph{difference}:
\[
\widehat{\text{Var}}(\widehat{\text{DDD}}) \;=\; \frac{1}{N}\,\widehat{\text{Var}}_n\!\left[\hat\psi^{\text{crit}}_i - \hat\psi^{\text{flat}}_i\right], \qquad N = \text{number of dates},
\]
i.e.\ the cross-class covariance is captured automatically by the joint distribution of IFs. The two-sided $p$-value is reported under asymptotic normality.
% Implementation: scripts/stata/reform_window/03_csdid_da15.do and 04_csdid_ida15.do

\textit{Why two estimators.} The OLS-FE DDD targets the standard triple-interaction coefficient $\beta_7$ under the unconditional PT assumption $\mathbb{E}[\Delta\Delta Y(0) \mid \text{y25}=1] = \mathbb{E}[\Delta\Delta Y(0) \mid \text{y25}=0]$. The DR-DiD targets the conditional ATT under the conditional PT $\mathbb{E}[Y_{\text{post}}(0) - Y_{\text{pre}}(0) \mid \text{y25}, X] = \mathbb{E}[Y_{\text{post}}(0) - Y_{\text{pre}}(0) \mid X]$, and is semiparametrically efficient under correct specification of the nuisance models (Sant'Anna--Zhao 2020). The two estimands coincide when $X$ is balanced across treatment cells; with calendar-matched windows, balance on $X = (\text{dow}, \text{month})$ is approximate but not exact (a year-on-year drift of $\approx 1$ day-of-week, plus daylight-saving boundaries). The discrepancy between the OLS-FE and DR-DiD point estimates is therefore the IPW contribution.

\paragraph{Naive before-after vs DDD: how much of the apparent reform effect is seasonal bias?}
The standard alternative to a DDD design is a within-year before-after (BA) regression with calendar (and sometimes weather) controls. Table~\ref{tab:ba_vs_ddd} compares the BA-with-calendar-controls estimate for each of our four (window, market) cells against the DDD with the same controls. The seasonal bias --- what the BA misattributes to the reform but the DDD nets out via the 2024 placebo --- is striking on the DA15 window: BA gives $+26$ EUR/MWh on the critical-flat spread, DDD gives essentially zero. The 2024 calendar-matched placebo year shows the same $+27$ swing in its own pre-vs-post comparison; almost the entire apparent BA effect is seasonal swing of the critical-flat spread from summer to winter, not the MTU15 reform. On the IDA15 window the bias is smaller (3--5 EUR/MWh) but goes the \emph{opposite direction} (the 2024 placebo shows a $\approx-8$ EUR/MWh swing, vs $+2$/$+6$ in the 2025 treatment year). In both windows the BA is uninformative about the reform; the DDD is the design that isolates the reform from the calendar.

\begin{table}[H]\centering\footnotesize
\setlength{\tabcolsep}{4pt}
\caption{Naive before-after (BA) vs DDD on the within-day (crit$-$flat) spread, EUR/MWh, all four (window, market) cells. BA: within-year (post $-$ pre) on 2025 spread, OLS with \code{i.dow + i.month} calendar controls and date-clustered SE. Placebo BA: same spec on 2024 spread. DDD: 2$\times$2 DiD with both years and same calendar controls (the headline OLS-FE coefficient on $\beta_{\text{post}\times\text{y25}}$). \emph{Seasonal bias} = BA $-$ DDD = what an advisor's naive within-year BA-with-controls attributes to the reform but is really seasonality (identified by the 2024 placebo). Standard errors in parentheses. Stars: $^{*}~0.10,~^{**}~0.05,~^{***}~0.01$.}
\label{tab:ba_vs_ddd}
\begin{tabular}{llcccc}
\toprule
Window & Outcome & BA 2025 & BA 2024 (placebo) & DDD & Seasonal bias \\
       &         & (a)     & (b)               & (c) & (a) $-$ (c)   \\
\midrule
DA15 & DA price &   25.83*** \scriptsize{( 3.72)} &   26.99*** \scriptsize{( 4.34)} &   -0.17 \scriptsize{( 2.41)} &   26.00 \\
DA15 & IDA-S1 price &   25.90*** \scriptsize{( 3.79)} &   26.42*** \scriptsize{( 3.99)} &   -2.24 \scriptsize{( 2.33)} &   28.15 \\
IDA15 & DA price &    2.24 \scriptsize{( 5.30)} &   -8.49*** \scriptsize{( 3.21)} &   -1.74 \scriptsize{( 3.35)} &    3.98 \\
IDA15 & IDA-S1 price &    6.07 \scriptsize{( 4.91)} &   -7.05** \scriptsize{( 3.25)} &    0.74 \scriptsize{( 3.31)} &    5.33 \\
\bottomrule
\end{tabular}
\end{table}


\paragraph{Headline DDD estimates: clearing prices, all (reform window, outcome) cells.}
Every price outcome is estimated in \emph{both} reform windows. The natural mapping (DA price tested in DA15 window, IDA price tested in IDA15 window) is the direct-effect cell; the off-diagonal cells (DA price in IDA15 window, IDA price in DA15 window) test \emph{spillover}. We report the doubly-robust DiD (Sant'Anna--Zhao 2020, \code{drdid}) on prices in levels, decomposed into ATT$_{\text{crit}}$ and ATT$_{\text{flat}}$, with the DDD given by the difference. Joint SE for the DDD comes from the difference of the recentered influence functions (drdid \code{stub()} option), which correctly handles the within-date covariance between the two hour-class means. Covariates: \code{i.dow + i.month}. Full per-cell OLS-FE companion tables: Tables~\ref{tab:reform_ddd_da15_da}--\ref{tab:reform_ddd_ida15_ida}.
% Pipeline: scripts/stata/reform_window/

\begin{table}[H]
\centering\footnotesize
\setlength{\tabcolsep}{4pt}
\caption{Headline DDD on \emph{prices in levels}, EUR/MWh, all four (reform window, market) cells. DR-DiD columns (Sant'Anna--Zhao 2020, drimp): ATT$_{\text{crit}}$ and ATT$_{\text{flat}}$ are 2$\times$2 DiDs per hour-class with covariates \code{i.dow + i.month}; DDD is their difference with joint SE from the influence-function difference. OLS-FE $\beta_7$ column: triple-interaction coefficient from \code{reghdfe} on the period-level panel with clock-hour + DOW + month FE, SE clustered at date. Standard errors in parentheses. Stars: $^{*}~0.10,~^{**}~0.05,~^{***}~0.01$. The \emph{Channel} column distinguishes direct (the market hit by the reform) from spillover (the cross-market test): DA15-DA and IDA15-IDA are direct; the other two cells are spillover.}
\label{tab:reform_ddd_summary}
\begin{tabular}{llcccc}
\toprule
& & \multicolumn{3}{c}{DR-DiD on levels} & OLS-FE \\
\cmidrule(lr){3-5}
Window & Outcome & ATT$_{\text{crit}}$ & ATT$_{\text{flat}}$ & DDD & $\beta_7$ \\
\midrule
DA15  & DA price        & $-8.34^{*}$ \scriptsize{(4.33)}    & $-2.96$ \scriptsize{(4.85)}     & $-5.38^{**}$ \scriptsize{(2.36)}  & $-0.29$ \scriptsize{(2.92)} \\
DA15  & IDA-S1 price    & $-10.38^{**}$ \scriptsize{(4.46)}  & $-5.02$ \scriptsize{(4.96)}     & $-5.36^{**}$ \scriptsize{(2.30)}  & $-2.67$ \scriptsize{(2.85)} \\
IDA15 & DA price        & $-25.02^{***}$ \scriptsize{(5.49)} & $-31.30^{***}$ \scriptsize{(6.16)} & $+6.28^{**}$ \scriptsize{(3.01)}  & $-2.24$ \scriptsize{(3.63)} \\
IDA15 & IDA-S1 price    & $-25.40^{***}$ \scriptsize{(5.59)} & $-34.82^{***}$ \scriptsize{(6.13)} & $+9.43^{***}$ \scriptsize{(2.95)} & $+0.50$ \scriptsize{(3.69)} \\
\bottomrule
\end{tabular}
\end{table}

\textit{OLS-FE vs DR-DiD.} Both estimators target the DDD parameter; the gap between them is the IPW reweighting on calendar covariates. Without covariates, DR-DiD reproduces OLS-FE to 3 decimal places (sanity check). With $X = (\text{i.dow}, \text{i.month})$, DR-DiD reweights observations so the (dow, month) distribution is balanced across 2024-placebo and 2025-treatment cells before computing the ATT. This balance is approximate but not exact under calendar-matched windows (day-of-week drifts $\approx 1$ position across years; DST boundaries differ by a day). The OLS-FE estimand is the unconditional uniform-weight DDD; the DR-DiD estimand is the conditional ATT under conditional PT. We report both. The DR-DiD is the preferred headline only insofar as conditional PT given $(\text{dow}, \text{month})$ is the assumption we are willing to maintain --- a discussion of covariate adequacy is below.

\textit{Reading the four price cells.}
DA15: both DA and IDA Session~1 spreads compress by $\approx 5.4$ EUR/MWh after MTU15-DA. Same direction, near-identical magnitude across markets ($-5.38$, $-5.36$). The IDA-cell is a clean spillover because MTU15-IDA is held constant within the DA15 window in both years. IDA15: both spreads widen ($+6.3$ DA, $+9.4$ IDA), again same direction across markets, with the directly-affected IDA larger. Hour-class decomposition: in DA15, critical and flat both fall ($-3$ to $-10$ EUR/MWh) with critical falling more; in IDA15, both fall sharply ($-25$ to $-35$, post-MTU15-IDA is structurally cheaper) with flat falling more. The DDD identifies the differential of these moves, not the absolute response of critical hours. Same-direction-within-window rules out a market-specific artefact; cross-window asymmetry (DA15 compresses, IDA15 widens) is a property of the reforms, not of the markets being measured.

\textit{Mechanism reading (pending joint $q_1$/$q_2$ check).}
The price DDDs are consistent with two compositional stories that we cannot separate from prices alone: (a) MTU15-DA opens 4$\times$ finer critical-hour quarter-slots that non-pivotal generators (renewables, hydro, fast-ramp gas) fill while pivotal firms reduce DA participation, compressing the spread; (b) MTU15-IDA reshapes IDA bidding so that the post-reform IDA market is structurally different in ways the DDD only captures via spread widening. The firm-level $q_1$/$q_2$ regressions below are the test that pins down which compositional story is consistent with the data.

\emph{Follow-ups flagged in priority order}: (a) BA-with-controls comparison to quantify how much of the apparent reform effect the DDD absorbs as seasonal; (b) closest-to-delivery IDA price (in place of Session 1) as a robustness for the IDA-spillover cell, given the June 2024 IDA reform redefined Session 1; (c) augmented DR-DiD covariate set with ENTSO-E wind+solar generation and TTF gas as exogenous structural cost shifters; (d) propensity-score overlap diagnostic to confirm DR-DiD weights are not extreme; (e) event-study with multiplier-bootstrap simultaneous CI.

\begin{table}[htbp]\centering
\def\sym#1{\ifmmode^{#1}\else\(^{#1}\)\fi}
\caption{OLS-FE DDD: DA clearing price in DA15 window. SE clustered at date.}\label{tab:reform_ddd_da15_da}
\begin{tabular}{l*{3}{cc}}
\toprule
                    &\multicolumn{2}{c}{(1)}     &\multicolumn{2}{c}{(2)}     &\multicolumn{2}{c}{(3)}     \\
                    &\multicolumn{2}{c}{Sparse}  &\multicolumn{2}{c}{Month FE}&\multicolumn{2}{c}{Date FE} \\
                    &           b   &          se&           b   &          se&           b   &          se\\
\midrule
crit*post*y25 ($\beta_7$, DDD)&      -0.286   &     (2.915)&      -0.286   &     (2.916)&      -0.286   &     (2.915)\\
crit*post           &      28.057***&     (2.073)&      28.057***&     (2.074)&      28.057***&     (2.073)\\
crit*y25            &      -3.529*  &     (1.951)&      -3.529*  &     (1.952)&      -3.529*  &     (1.951)\\
post*y25            &     -26.523***&     (5.983)&     -26.523***&     (5.335)&               &            \\
crit                &       0.000   &         (.)&       0.000   &         (.)&       0.000   &         (.)\\
post                &      11.846***&     (4.112)&       0.000   &         (.)&               &            \\
y25                 &      -1.662   &     (4.010)&      -1.662   &     (2.947)&               &            \\
\midrule
Observations        &       13888   &            &       13888   &            &       13888   &            \\
$R^2$               &       0.282   &            &       0.416   &            &       0.800   &            \\
Clock-hour FE       &         Yes   &            &         Yes   &            &         Yes   &            \\
DOW FE              &         Yes   &            &         Yes   &            &         Yes   &            \\
Month FE            &          No   &            &         Yes   &            &         n/a   &            \\
Date FE             &          No   &            &          No   &            &         Yes   &            \\
Cluster             &        Date   &            &        Date   &            &        Date   &            \\
\bottomrule
\end{tabular}
\end{table}

\begin{table}[htbp]\centering
\def\sym#1{\ifmmode^{#1}\else\(^{#1}\)\fi}
\caption{OLS-FE DDD: IDA Session 1 clearing price in DA15 window. SE clustered at date.}\label{tab:reform_ddd_da15_ida}
\begin{tabular}{l*{3}{cc}}
\toprule
                    &\multicolumn{2}{c}{(1)}     &\multicolumn{2}{c}{(2)}     &\multicolumn{2}{c}{(3)}     \\
                    &\multicolumn{2}{c}{Sparse}  &\multicolumn{2}{c}{Month FE}&\multicolumn{2}{c}{Date FE} \\
                    &           b   &          se&           b   &          se&           b   &          se\\
\midrule
crit*post*y25 ($\beta_7$, DDD)&      -2.674   &     (2.852)&      -2.674   &     (2.853)&      -2.674   &     (2.852)\\
crit*post           &      28.703***&     (1.987)&      28.703***&     (1.988)&      28.703***&     (1.987)\\
crit*y25            &      -2.421   &     (1.958)&      -2.421   &     (1.958)&      -2.421   &     (1.958)\\
post*y25            &     -27.432***&     (6.030)&     -26.886***&     (5.313)&               &            \\
crit                &       0.000   &         (.)&       0.000   &         (.)&       0.000   &         (.)\\
post                &      11.870***&     (4.038)&       0.000   &         (.)&               &            \\
y25                 &      -3.158   &     (4.075)&      -3.210   &     (2.961)&               &            \\
\midrule
Observations        &       19138   &            &       19138   &            &       19138   &            \\
$R^2$               &       0.299   &            &       0.479   &            &       0.785   &            \\
Clock-hour FE       &         Yes   &            &         Yes   &            &         Yes   &            \\
DOW FE              &         Yes   &            &         Yes   &            &         Yes   &            \\
Month FE            &          No   &            &         Yes   &            &         n/a   &            \\
Date FE             &          No   &            &          No   &            &         Yes   &            \\
Cluster             &        Date   &            &        Date   &            &        Date   &            \\
\bottomrule
\end{tabular}
\end{table}

\begin{table}[htbp]\centering
\def\sym#1{\ifmmode^{#1}\else\(^{#1}\)\fi}
\caption{OLS-FE DDD: DA clearing price in IDA15 window. SE clustered at date.}\label{tab:reform_ddd_ida15_da}
\begin{tabular}{l*{3}{cc}}
\toprule
                    &\multicolumn{2}{c}{(1)}     &\multicolumn{2}{c}{(2)}     &\multicolumn{2}{c}{(3)}     \\
                    &\multicolumn{2}{c}{Sparse}  &\multicolumn{2}{c}{Month FE}&\multicolumn{2}{c}{Date FE} \\
                    &           b   &          se&           b   &          se&           b   &          se\\
\midrule
crit*post*y25 ($\beta_7$, DDD)&      -2.243   &     (3.630)&      -2.243   &     (3.631)&      -2.243   &     (3.628)\\
crit*post           &     -10.440***&     (1.738)&     -10.440***&     (1.739)&     -10.440***&     (1.738)\\
crit*y25            &       2.657   &     (2.282)&       2.657   &     (2.284)&       2.657   &     (2.282)\\
post*y25            &     -27.426***&     (6.459)&     -27.438***&     (6.250)&               &            \\
crit                &       0.000   &         (.)&       0.000   &         (.)&       0.000   &         (.)\\
post                &     -34.767***&     (4.266)&      -3.141   &     (6.004)&               &            \\
y25                 &      45.676***&     (4.598)&      45.670***&     (4.229)&               &            \\
\midrule
Observations        &        3934   &            &        3934   &            &        3934   &            \\
$R^2$               &       0.573   &            &       0.627   &            &       0.874   &            \\
Clock-hour FE       &         Yes   &            &         Yes   &            &         Yes   &            \\
DOW FE              &         Yes   &            &         Yes   &            &         Yes   &            \\
Month FE            &          No   &            &         Yes   &            &         n/a   &            \\
Date FE             &          No   &            &          No   &            &         Yes   &            \\
Cluster             &        Date   &            &        Date   &            &        Date   &            \\
\bottomrule
\end{tabular}
\end{table}

\begin{table}[htbp]\centering
\def\sym#1{\ifmmode^{#1}\else\(^{#1}\)\fi}
\caption{OLS-FE DDD: IDA Session 1 clearing price in IDA15 window. SE clustered at date.}\label{tab:reform_ddd_ida15_ida}
\begin{tabular}{l*{3}{cc}}
\toprule
                    &\multicolumn{2}{c}{(1)}     &\multicolumn{2}{c}{(2)}     &\multicolumn{2}{c}{(3)}     \\
                    &\multicolumn{2}{c}{Sparse}  &\multicolumn{2}{c}{Month FE}&\multicolumn{2}{c}{Date FE} \\
                    &           b   &          se&           b   &          se&           b   &          se\\
\midrule
crit*post*y25 ($\beta_7$, DDD)&       0.504   &     (3.694)&       0.504   &     (3.695)&       0.504   &     (3.693)\\
crit*post           &      -9.933***&     (1.693)&      -9.933***&     (1.694)&      -9.933***&     (1.693)\\
crit*y25            &       3.266   &     (2.155)&       3.266   &     (2.155)&       3.266   &     (2.154)\\
post*y25            &     -29.831***&     (6.566)&     -29.975***&     (6.342)&               &            \\
crit                &       0.000   &         (.)&       0.000   &         (.)&       0.000   &         (.)\\
post                &     -36.304***&     (4.397)&      -6.630   &     (5.983)&               &            \\
y25                 &      45.197***&     (4.585)&      45.316***&     (4.202)&               &            \\
\midrule
Observations        &        5530   &            &        5530   &            &        5530   &            \\
$R^2$               &       0.588   &            &       0.627   &            &       0.837   &            \\
Clock-hour FE       &         Yes   &            &         Yes   &            &         Yes   &            \\
DOW FE              &         Yes   &            &         Yes   &            &         Yes   &            \\
Month FE            &          No   &            &         Yes   &            &         n/a   &            \\
Date FE             &          No   &            &          No   &            &         Yes   &            \\
Cluster             &        Date   &            &        Date   &            &        Date   &            \\
\bottomrule
\end{tabular}
\end{table}

\begin{table}[H]\centering\small
\caption{DR-DiD on DA clearing price in DA15 window, prices in levels (Sant'Anna--Zhao 2020). 2$\times$2 DiD per hour-class, covariates \texttt{i.dow + i.month}; DDD is the difference of ATTs with joint SE from the influence-function difference.}\label{tab:reform_csdid_da15_da}
\begin{tabular}{lrr}
\toprule
Coefficient & Estimate & SE \\\\
\midrule
ATT$_{\text{crit}}$ &  -8.336 & (4.326) \\\\
ATT$_{\text{flat}}$ &  -2.961 & (4.846) \\\\
DDD = ATT$_{\text{crit}}$ $-$ ATT$_{\text{flat}}$ &  -5.375 & (2.356) \\\\
\midrule
p-value (DDD)        & 0.023 & \\\\
95\% CI (DDD)        & [ -9.992,  -0.758] & \\\\
N (dates)            &   584 & \\\\
Outcome              & DA clearing price & \\\\
Window               & DA15 (reform = 2025-10-01) & \\\\
\bottomrule
\end{tabular}
\end{table}

\begin{table}[H]\centering\small
\caption{DR-DiD on IDA Session 1 clearing price in DA15 window, prices in levels (Sant'Anna--Zhao 2020). 2$\times$2 DiD per hour-class, covariates \texttt{i.dow + i.month}; DDD is the difference of ATTs with joint SE from the influence-function difference.}\label{tab:reform_csdid_da15_ida}
\begin{tabular}{lrr}
\toprule
Coefficient & Estimate & SE \\\\
\midrule
ATT$_{\text{crit}}$ & -10.379 & (4.463) \\\\
ATT$_{\text{flat}}$ &  -5.016 & (4.957) \\\\
DDD = ATT$_{\text{crit}}$ $-$ ATT$_{\text{flat}}$ &  -5.364 & (2.296) \\\\
\midrule
p-value (DDD)        & 0.019 & \\\\
95\% CI (DDD)        & [ -9.864,  -0.863] & \\\\
N (dates)            &   554 & \\\\
Outcome              & IDA Session 1 clearing price & \\\\
Window               & DA15 (reform = 2025-10-01) & \\\\
\bottomrule
\end{tabular}
\end{table}

\begin{table}[H]\centering\small
\caption{DR-DiD on DA clearing price in IDA15 window, prices in levels (Sant'Anna--Zhao 2020). 2$\times$2 DiD per hour-class, covariates \texttt{i.dow + i.month}; DDD is the difference of ATTs with joint SE from the influence-function difference.}\label{tab:reform_csdid_ida15_da}
\begin{tabular}{lrr}
\toprule
Coefficient & Estimate & SE \\\\
\midrule
ATT$_{\text{crit}}$ & -25.019 & (5.486) \\\\
ATT$_{\text{flat}}$ & -31.302 & (6.161) \\\\
DDD = ATT$_{\text{crit}}$ $-$ ATT$_{\text{flat}}$ &   6.283 & (3.009) \\\\
\midrule
p-value (DDD)        & 0.037 & \\\\
95\% CI (DDD)        & [  0.386,  12.180] & \\\\
N (dates)            &   281 & \\\\
Outcome              & DA clearing price & \\\\
Window               & IDA15 (reform = 2025-03-19) & \\\\
\bottomrule
\end{tabular}
\end{table}

\begin{table}[H]\centering\small
\caption{DR-DiD on IDA Session 1 clearing price in IDA15 window, prices in levels (Sant'Anna--Zhao 2020). 2$\times$2 DiD per hour-class, covariates \texttt{i.dow + i.month}; DDD is the difference of ATTs with joint SE from the influence-function difference.}\label{tab:reform_csdid_ida15_ida}
\begin{tabular}{lrr}
\toprule
Coefficient & Estimate & SE \\\\
\midrule
ATT$_{\text{crit}}$ & -25.396 & (5.587) \\\\
ATT$_{\text{flat}}$ & -34.823 & (6.129) \\\\
DDD = ATT$_{\text{crit}}$ $-$ ATT$_{\text{flat}}$ &   9.427 & (2.949) \\\\
\midrule
p-value (DDD)        & 0.001 & \\\\
95\% CI (DDD)        & [  3.647,  15.207] & \\\\
N (dates)            &   278 & \\\\
Outcome              & IDA Session 1 clearing price & \\\\
Window               & IDA15 (reform = 2025-03-19) & \\\\
\bottomrule
\end{tabular}
\end{table}


\paragraph{IDA-price robustness: closest-to-delivery session.}
The IDA-cell DDDs in Table~\ref{tab:reform_ddd_summary} use Session 1 price throughout. Session 1's clearing horizon was redefined at the SIDC reform of June 2024 (MIBEL 6-session $\to$ SIDC 3-session, per \code{docs/notes/SPANISH_MARKET_STRUCTURE.md} \S IDA): under MIBEL Session 1 closed at 14:00 D$-$1 and targeted all 24 hours of D; under SIDC it covers all 96 ISP15 periods of D. Both target the same horizon nominally, but liquidity, participants, and the SIDC coupling regime differ in ways the within-window pre/post differencing cannot absorb when the placebo year is MIBEL and the treatment year is SIDC. As a robustness, we re-estimate the IDA-cell DR-DiD using the \emph{closest-to-delivery} session price per (date, period) --- the row with the maximum \code{session_number}, which under SIDC is IDA-3 for afternoons and IDA-2 for mornings, and under MIBEL is Session 4--6 for D-day afternoons and Session 3 for mornings. This is MTU- and session-count-invariant by construction. Results in Table~\ref{tab:csdid_ida_latest_compare}:

\begin{table}[H]\centering\footnotesize
\setlength{\tabcolsep}{4pt}
\caption{DR-DiD on IDA prices: Session 1 (baseline, in Table~\ref{tab:reform_ddd_summary}) vs closest-to-delivery session (robustness). For each (date, period) the closest-to-delivery price is the row with the maximum \code{session_number}: under SIDC, IDA-3 for afternoon and IDA-2 for morning; under MIBEL, Sessions 4--6 for D-day afternoons and Session 3 for morning. This robustness removes the MIBEL-to-SIDC redefinition of Session 1. Covariates: \code{i.dow + i.month}. Standard errors in parentheses. Stars: $^{*}$ 0.10, $^{**}$ 0.05, $^{***}$ 0.01.}
\label{tab:csdid_ida_latest_compare}
\begin{tabular}{lccc}
\toprule
Window & ATT$_{\text{crit}}$ & ATT$_{\text{flat}}$ & DDD \\
\midrule
DA15 &   -2.54 \scriptsize{( 4.39)} &   -3.76 \scriptsize{( 4.75)} &    1.21 \scriptsize{( 2.33)} \\
IDA15 &  -27.84 \scriptsize{( 5.16)} &  -34.10 \scriptsize{( 5.88)} &    6.26* \scriptsize{( 3.38)} \\
\bottomrule
\end{tabular}
\end{table}


The DA15-IDA "spillover" cell --- baseline DDD of $-5.36$ EUR/MWh (significant) --- effectively disappears under the robustness ($+1.21$, $p=0.60$). The IDA15-IDA "direct" cell shrinks from $+9.43^{***}$ to $+6.26$ ($p=0.064$), losing conventional significance. \textbf{The two IDA-side DDD numbers in the headline table are partly artefacts of the Session 1 choice}: the DA15 spillover is not robust at all, and the IDA15 direct effect is robust in direction and order of magnitude but loses statistical significance. The DA-side cells of the headline table (DA15-DA and IDA15-DA) are unaffected by this robustness (no session choice in the DA market). We will treat the IDA-side numbers in the headline as upper bounds on the true DDD until the wild-cluster-bootstrap and closest-to-delivery results converge.

\paragraph{Propensity-score overlap diagnostic.}
The DR-DiD propensity score $\hat p(X) = \Pr(\text{y25} = 1 \mid \text{i.dow}, \text{i.month})$ is estimated by logit on each (window, market) spread panel; weights blow up if $\hat p(X)$ clusters near 0 or 1. Across all four cells, $\hat p(X)$ ranges $[0.38, 0.54]$ (DA15-IDA) to $[0.49, 0.51]$ (DA15-DA), with median $\approx 0.50$ in every cell and 0\% of observations below 0.10 or above 0.90. Calendar-matched windows deliver essentially balanced treatment assignment in (dow, month) cells by design, so the IPW weights are not extreme and the DR-DiD inference is not driven by propensity-score tail behaviour. Histograms at \code{figures/working/fig_overlap_{window}_{market}.pdf}.

\paragraph{Headline DDD estimates: firm-level outcomes ($q_1$, $q_2$).}
Outcomes at the firm level: $q_1$ = DA-cleared MWh per (firm, unit, date, clock-hour) from OMIE \code{pdbc}, $q_2$ = signed intraday programmed MWh from \code{pibci}, both zero-filled for offline cells. CCGT-only sample: 42 pivotal Big-4 units (IB, GE, GN, HC, EDP-PT) and 8 non-pivotal units (Repsol, Engie, TotalEnergies, Moeve); untagged firms dropped. We report the pooled DDD ($\beta_7$) and the pivotal-status heterogeneity contrast $\delta_{\text{piv}} = \widehat{\text{DDD}}_{\text{pivotal}} - \widehat{\text{DDD}}_{\text{non-pivotal}}$ (per the framing above; this is what we previously over-labelled $\beta_{1234}$). All specs include firm + unit + clock-hour + DOW fixed effects. Full per-cell tables: Tables~\ref{tab:reform_ddd_firm_q1_da15}--\ref{tab:reform_ddd_firm_q2_ida15}.

\begin{table}[H]
\centering\footnotesize
\setlength{\tabcolsep}{4pt}
\caption{Pooled DDD ($\beta_7$) and pivotal-status heterogeneity contrast ($\delta_{\text{piv}}$) on firm-level $q_1$ and $q_2$, MWh per unit-hour. $\delta_{\text{piv}} > 0$ means pivotal firms' DDD effect is \emph{larger} (in algebraic terms) than non-pivotal firms'. SE clustered at date (pooled) and at firm (decomposed; $G \approx 9$, asymptotic SE may under-cover, wild-cluster bootstrap recommended). Standard errors in parentheses. Stars: $^{*}~0.10,~^{**}~0.05,~^{***}~0.01$.}
\label{tab:reform_firm_summary}
\begin{tabular}{llcc|cc}
\toprule
& & \multicolumn{2}{c}{Pooled DDD ($\beta_7$)} & \multicolumn{2}{c}{$\delta_{\text{piv}}$ (heterogeneity)} \\
\cmidrule(lr){3-4}\cmidrule(lr){5-6}
Outcome & Window & $\beta_7$ & SE & $\delta_{\text{piv}}$ & SE \\
\midrule
$q_1$ & DA15  & $-16.27^{***}$ & (2.69) & $-18.65^{**}$  & (7.35) \\
$q_1$ & IDA15 & $-9.96^{***}$  & (3.10) & $-4.01$        & (6.55) \\
$q_2$ & DA15  & $-1.27^{**}$   & (0.64) & $+18.78$       & (15.82) \\
$q_2$ & IDA15 & $-1.70^{**}$   & (0.85) & $+11.10$       & (10.81) \\
\bottomrule
\end{tabular}
\end{table}

\textit{Joint $(q_1, q_2)$ reading.}
Locating each window's pooled $\beta_7$ on the Table~\ref{tab:reform_outcomes} sign-matrix: DA15 lands in $(\beta(q_1) < 0,\,\beta(q_2) < 0)$ and IDA15 lands in the same quadrant. Both reforms produce \textbf{broader disengagement}, not Ito--Reguant withholding (which would require $\beta(q_2) > 0$): dispatchable CCGT firms reduce both DA-cleared volume and intraday programmed volume. The natural reading is that the relevant outside option for these units is not the IDA leg but REE's post-clearing redispatch (RT2 pay-as-bid; characterised in the descriptive companion, Part B). This is a substantively different story from the classical Ito--Reguant signature.

\textit{Pivotal-status heterogeneity reading.}
(i) For $q_1$ in the DA15 window the pooled DDD effect is sharply concentrated in pivotal firms: $\delta_{\text{piv}} = -18.65$ MWh/unit-hour, with the pooled $\beta_7$ ($-16.27$) entirely absorbed by the heterogeneity contrast. Non-pivotal firms in this window show no measurable change. (ii) For $q_1$ in the IDA15 window the pooled effect is significant ($-9.96$) but $\delta_{\text{piv}}$ is small and non-significant: the IDA reform pulled all CCGT firms in the same direction, no market-power moderation. (iii) For $q_2$, $\delta_{\text{piv}}$ has wide SEs and is not identifiable in either window. (iv) These contrasts identify heterogeneity \emph{in the DDD effect}, not the absolute effect on pivotal firms --- they tell us where the reform's response is concentrated, not whether the response is "strategic" in a structural sense.

\textit{Tension with the older 2-way DiD in \code{paper.tex}.}
The appendix in \code{paper.tex} (commented-out main table) reports a 2-way DiD (no 2024 placebo) on $q_2$ in pivotal vs non-pivotal subsamples with \emph{opposite signs} (+3.87 pivotal vs $-$4.74 non-pivotal, both highly significant). Our DDD on the same outcome shows both subsamples moving in the \emph{same} direction (both negative). The sign reversal in the older 2-way DiD was driven by the seasonality that the 2024 placebo here absorbs. Reconciling: once we net out YoY seasonality, the pivotal-vs-non-pivotal contrast in $q_2$ is wide-SE and uninformative --- the headline pattern in the appendix was partly seasonal.

\textit{Caveats.}
(a) Firm-cluster $G \approx 9$ is small; asymptotic SE may under-cover. Wild-cluster bootstrap (\code{boottest}) is a planned robustness. (b) Pivotal/non-pivotal partition treats all 5 Big-4 firms as identical and all 4 placebo firms as identical; firm-by-firm decomposition is a natural next step (Moeve has a documented pre-trend break in \code{paper.tex} that may pull the non-pivotal subsample). (c) The 2024 placebo for the DA15 $q_1$ panel includes 6 weeks (Apr 28 -- Jun 13, 2024) before the June 14, 2024 IDA reform changed sessions $6 \to 3$. Robustness fix: use closest-to-delivery IDA price (MTU- and session-count-invariant) instead of Session 1 for the IDA-spillover regressions.

\begin{table}[htbp]\centering
\def\sym#1{\ifmmode^{#1}\else\(^{#1}\)\fi}
\caption{DDD and DDDD on $q\_1$ (DA-cleared MWh/unit-hour), da15 window.}\label{tab:reform_ddd_firm_q1_da15}
\begin{tabular}{l*{2}{cc}}
\toprule
                    &\multicolumn{2}{c}{(1)}     &\multicolumn{2}{c}{(2)}     \\
                    &\multicolumn{2}{c}{DDD}     &\multicolumn{2}{c}{DDDD}    \\
                    &           b   &          se&           b   &          se\\
\midrule
crit*post*y25*piv ($\beta_{1234}$, DDDD)&               &            &      -18.65** &      (7.35)\\
crit*post*y25 ($\beta_7$, DDD)&      -16.27***&      (2.69)&       -0.60   &      (1.77)\\
crit*post*piv       &               &            &       13.61   &     (11.04)\\
crit*y25*piv        &               &            &        4.37   &      (7.82)\\
post*y25*piv        &               &            &       18.96   &     (30.23)\\
crit*piv            &               &            &       -0.07   &      (2.28)\\
post*piv            &               &            &      -11.81   &     (10.44)\\
y25*piv             &               &            &       -4.40   &     (10.81)\\
crit*post           &       22.37***&      (2.43)&       10.95   &      (9.18)\\
crit*y25            &       -2.00   &      (1.40)&       -5.68   &      (7.15)\\
post*y25            &      -12.01***&      (3.37)&      -27.93   &     (29.98)\\
crit                &        0.00   &         (.)&        0.00   &         (.)\\
post                &       -3.14   &      (2.47)&        6.77   &      (9.44)\\
y25                 &        3.98   &      (2.77)&        7.68   &      (9.75)\\
\midrule
Observations        &      408800   &            &      408800   &            \\
$R^2$               &       0.092   &            &       0.092   &            \\
Spec                &         DDD   &            &        DDDD   &            \\
Cluster             &        Date   &            &        Firm   &            \\
Unit FE             &         Yes   &            &         Yes   &            \\
Clock-hour FE       &         Yes   &            &         Yes   &            \\
DOW FE              &         Yes   &            &         Yes   &            \\
\bottomrule
\end{tabular}
\end{table}

\begin{table}[htbp]\centering
\def\sym#1{\ifmmode^{#1}\else\(^{#1}\)\fi}
\caption{DDD and DDDD on $q\_1$ (DA-cleared MWh/unit-hour), ida15 window.}\label{tab:reform_ddd_firm_q1_ida15}
\begin{tabular}{l*{2}{cc}}
\toprule
                    &\multicolumn{2}{c}{(1)}     &\multicolumn{2}{c}{(2)}     \\
                    &\multicolumn{2}{c}{DDD}     &\multicolumn{2}{c}{DDDD}    \\
                    &           b   &          se&           b   &          se\\
\midrule
crit*post*y25*piv ($\beta_{1234}$, DDDD)&               &            &       -4.01   &      (6.55)\\
crit*post*y25 ($\beta_7$, DDD)&       -9.96***&      (3.10)&       -6.60   &      (6.16)\\
crit*post*piv       &               &            &       -6.86*  &      (3.07)\\
crit*y25*piv        &               &            &        3.83   &      (6.54)\\
post*y25*piv        &               &            &       14.64   &     (17.72)\\
crit*piv            &               &            &        8.01*  &      (3.58)\\
post*piv            &               &            &       -2.88   &      (2.39)\\
y25*piv             &               &            &      -12.55   &     (17.89)\\
crit*post           &       -6.95***&      (1.61)&       -1.18*  &      (0.56)\\
crit*y25            &       10.00***&      (3.05)&        6.79   &      (6.12)\\
post*y25            &       -7.18***&      (2.49)&      -19.48   &     (17.61)\\
crit                &        0.00   &         (.)&        0.00   &         (.)\\
post                &       -1.11   &      (0.94)&        1.31   &      (2.35)\\
y25                 &        7.10***&      (2.28)&       17.65   &     (17.82)\\
\midrule
Observations        &      196700   &            &      196700   &            \\
$R^2$               &       0.072   &            &       0.073   &            \\
Spec                &         DDD   &            &        DDDD   &            \\
Cluster             &        Date   &            &        Firm   &            \\
Unit FE             &         Yes   &            &         Yes   &            \\
Clock-hour FE       &         Yes   &            &         Yes   &            \\
DOW FE              &         Yes   &            &         Yes   &            \\
\bottomrule
\end{tabular}
\end{table}

\begin{table}[htbp]\centering
\def\sym#1{\ifmmode^{#1}\else\(^{#1}\)\fi}
\caption{DDD and DDDD on $q\_2$ (intraday MWh/unit-hour), da15 window.}\label{tab:reform_ddd_firm_q2_da15}
\begin{tabular}{l*{2}{cc}}
\toprule
                    &\multicolumn{2}{c}{(1)}     &\multicolumn{2}{c}{(2)}     \\
                    &\multicolumn{2}{c}{DDD}     &\multicolumn{2}{c}{DDDD}    \\
                    &           b   &          se&           b   &          se\\
\midrule
crit*post*y25*piv ($\beta_{1234}$, DDDD)&               &            &       18.78   &     (15.82)\\
crit*post*y25 ($\beta_7$, DDD)&       -1.27** &      (0.64)&      -17.04   &     (15.78)\\
crit*post*piv       &               &            &      -19.42   &     (15.61)\\
crit*y25*piv        &               &            &       21.45   &     (19.14)\\
post*y25*piv        &               &            &       -2.77   &      (2.42)\\
crit*piv            &               &            &      -21.90   &     (19.30)\\
post*piv            &               &            &        2.61   &      (2.55)\\
y25*piv             &               &            &       -4.20***&      (1.11)\\
crit*post           &        1.79***&      (0.62)&       18.10   &     (15.57)\\
crit*y25            &       -2.11***&      (0.40)&      -20.12   &     (19.14)\\
post*y25            &       -2.96***&      (0.71)&       -0.63   &      (0.76)\\
crit                &        0.00   &         (.)&        0.00   &         (.)\\
post                &        2.51***&      (0.69)&        0.32   &      (0.91)\\
y25                 &       -4.30***&      (0.47)&       -0.77   &      (0.78)\\
\midrule
Observations        &      408800   &            &      408800   &            \\
$R^2$               &       0.105   &            &       0.133   &            \\
Spec                &         DDD   &            &        DDDD   &            \\
Cluster             &        Date   &            &        Firm   &            \\
Unit FE             &         Yes   &            &         Yes   &            \\
Clock-hour FE       &         Yes   &            &         Yes   &            \\
DOW FE              &         Yes   &            &         Yes   &            \\
\bottomrule
\end{tabular}
\end{table}

\begin{table}[htbp]\centering
\def\sym#1{\ifmmode^{#1}\else\(^{#1}\)\fi}
\caption{DDD and DDDD on $q\_2$ (intraday MWh/unit-hour), ida15 window.}\label{tab:reform_ddd_firm_q2_ida15}
\begin{tabular}{l*{2}{cc}}
\toprule
                    &\multicolumn{2}{c}{(1)}     &\multicolumn{2}{c}{(2)}     \\
                    &\multicolumn{2}{c}{DDD}     &\multicolumn{2}{c}{DDDD}    \\
                    &           b   &          se&           b   &          se\\
\midrule
crit*post*y25*piv ($\beta_{1234}$, DDDD)&               &            &       11.10   &     (10.81)\\
crit*post*y25 ($\beta_7$, DDD)&       -1.70** &      (0.85)&      -11.03   &     (10.19)\\
crit*post*piv       &               &            &       24.79   &     (21.64)\\
crit*y25*piv        &               &            &       -7.27   &      (8.63)\\
post*y25*piv        &               &            &        1.46   &      (4.09)\\
crit*piv            &               &            &      -30.97   &     (26.45)\\
post*piv            &               &            &       -5.40*  &      (2.85)\\
y25*piv             &               &            &       -3.39   &      (4.50)\\
crit*post           &       -1.81***&      (0.63)&      -22.63   &     (21.42)\\
crit*y25            &        0.38   &      (0.73)&        6.49   &      (7.70)\\
post*y25            &       -1.95** &      (0.80)&       -3.17   &      (2.16)\\
crit                &        0.00   &         (.)&        0.00   &         (.)\\
post                &       -4.77***&      (0.60)&       -0.24   &      (0.15)\\
y25                 &        0.55   &      (0.65)&        3.40   &      (2.34)\\
\midrule
Observations        &      196700   &            &      196700   &            \\
$R^2$               &       0.162   &            &       0.184   &            \\
Spec                &         DDD   &            &        DDDD   &            \\
Cluster             &        Date   &            &        Firm   &            \\
Unit FE             &         Yes   &            &         Yes   &            \\
Clock-hour FE       &         Yes   &            &         Yes   &            \\
DOW FE              &         Yes   &            &         Yes   &            \\
\bottomrule
\end{tabular}
\end{table}


\paragraph{Granularity strategy.}
The cross-reform DDD runs on disaggregated period-level data, not on hourly aggregates. Pre-MTU15 the period is a 60-minute slot (1 obs / hour); post-MTU15 the period is a 15-minute slot (4 obs / hour). Both are valid observations of the unit's strategic choice within the relevant clock-hour, so we stack them in the same regression and let the FE structure absorb the time-aggregation difference. The \code{crit}, \code{post}, \code{y25} indicators are defined at the clock-hour level and replicate across the 4 quarters in the post-period. This gives $4\times$ the observations per hour post-reform, yielding tighter SE; SE are clustered at the date level to preserve correct inference under within-day dependence.

The within-hour quarter index is not used as a treatment-defining variable: ex ante there is no reason firms should bid more aggressively in a particular within-hour quarter regardless of clock-hour (e.g., they do not systematically bid less aggressively in the first quarter of every hour). The granularity is exploited as residual heterogeneity that the FE absorb, not as a treatment dimension. Post-reform cross-section descriptive use of the 15-min data (the $D_w$, $\Delta p$, $\Delta q$ histograms in the descriptive companion) is a separate complementary exercise.

\paragraph{Supporting evidence elsewhere in the project.}
See the sibling document \code{thesis/provisional/descriptive_facts.tex} for: (i) RT2 channel characterisation (Part B: REE post-clearing CCGT inclusion, Path A: $+1\,044$ GWh/mo CCGT post-blackout); (ii) zonal concentration (Part B: GN-dominated Cataluña/Levante/Sur, GE-Galicia, IB-Centro); (iii) GN's 840 EUR/MWh bid plateau invariant across regime (Part B: DA bid prices on RT2-up days); (iv) CNMC 2019 sanction (Part B: Cournot reading + CNMC); (v) $q_1$ collapse same-cal-month CCGT-only with parallel-trends on the post-clearing gap (Part A: $q_1$ drop).

% \paragraph{Source data.}
% \code{scripts/analysis/firm/reform_windows_timeline.py} (timeline);
% \code{scripts/analysis/firm/reform_window_event_study.py} (reduced-form event-study panel and figure);
% \code{scripts/stata/reform_window/} (headline DDD + DR-DiD robustness pipeline);
% panels at \code{results/regressions/firm/reform_window/panels/}; coefficients at \code{results/regressions/firm/reform_window/coefs/}; LaTeX tables at \code{thesis/paper/tables/reform_window/}.

\end{document}
